\documentclass[12pt]{article}
\usepackage[utf8]{inputenc}
\usepackage[T1]{fontenc}
\usepackage{hyperref}
\title{Porohanon Language Resources for K-12 Curriculum Development}
\author{Compiled by Claude Sonnet 5.5 (claude-sonnet-5-5) for Andr\'as Kornai}
\date{}
\begin{document}
\maketitle
\section{Introduction}
Porohanon is a Bisayan language spoken in the Camotes Islands, Cebu, Philippines. Web searches found only a small body of academic material and no school materials, primers, dictionaries or audio collections in Porohanon. This report lists only what appeared in search results.
\section{Primary Educational Resources}
No primers, textbooks or school programmes for Porohanon were found in the searches.
\section{Linguistic and Reference Materials}
\begin{itemize}
\item A Grammatical Sketch of Porohanon (V.\ C.\ Santiago, 2023 thesis, University of the Philippines Diliman): \url{https://linguistics.upd.edu.ph/theses_dissertation/a-grammatical-sketch-of-porohanon}
\item Porohanon Grammar Addenda (Santiago, 2024 presentation): \url{https://pages.upd.edu.ph/vasantiago/presentations/porohanon-grammar-addenda}
\item Language and Place in Poro, Camotes, Cebu (Santiago): \url{https://pages.upd.edu.ph/vasantiago/node/4447}
\item Porohanon language (Wikipedia overview): \url{https://en.wikipedia.org/wiki/Porohanon_language}
\end{itemize}
\section{Religious and Community Resources}
None found.
\section{Digital and Audio Resources}
None found; no dictionary or wordlist appeared in the results.
\section{Curriculum Development Context}
The language is closely related to Hiligaynon and shares much vocabulary with Cebuano, which heavily influences it. A curriculum would have to start from the 2023 grammar sketch and new fieldwork, since no pedagogical materials were found.
\end{document}